% !TEX root = ../../../main.tex
% 所属：第一部分：单变量分析学 / 第4章：数列极限的应用与一些反例
% 节：4.1 数列极限的定义与性质的应用
% 正文按原文件顺序保存；由 ../chapter4.tex 调用。
\section{数列极限的定义与性质的应用}
\begin{example}
	用极限定义证明 $\lim\limits _{n \rightarrow \infty} \dfrac{n}{n+3}=1$
	
\end{example}
\begin{proof}
	$\forall \varepsilon>0$ ，考虑
	
	$$
	\begin{aligned}
		& \left|\frac{n}{n+3}-1\right|<\varepsilon \\
		& \left|\frac{n}{n+3}-1\right|=\left|\frac{3}{n+3}\right|<\left|\frac{3}{n}\right|<\varepsilon
	\end{aligned}
	$$
	
	
	取 $N=\left[\dfrac{3}{\varepsilon}\right]+10$
	当 $n>N$ 时，必有 $\left|\dfrac{n}{n-3}-1\right|<\left|\dfrac{3}{n}\right|<\varepsilon$
\end{proof}
\begin{example}
	证明 $\left\{q^n\right\}, q \in(0,1)$ 是无穷小量
\end{example}
\begin{proof}
	$\forall \varepsilon>0$ ，考虑
	$$	\left|q^n\right|<\varepsilon $$
	取对数
	$$n \lg q<\lg \varepsilon$$
	
	
	取 $N=\left[\dfrac{\lg \varepsilon}{\lg q}\right]+10$
	
	$\forall n>N$ ，必有 $\left|q^n\right|<\varepsilon$
\end{proof}
\begin{proposition}\label{几个特殊的序列}
	几个特殊的序列
	\begin{enumerate}
		\item $a>0$ 时 $~\lim\limits _{n \rightarrow \infty} \dfrac{1}{n^a}=0$
		\item $a>0$时 $~\lim\limits _{n \rightarrow \infty} \sqrt[n]{a}=0$
		\item $\lim\limits _{n \rightarrow \infty} \sqrt[n]{n}=1$
		更一般地 $\forall k \in \mathbb{Z}_{+}$
		
		$$
		\lim _{n \rightarrow \infty} \sqrt[n]{n^k}=1
		$$
		\item $p>0$ ,而 $\alpha \in \mathbb{R}$
		
		$$
		\lim _{n \rightarrow \infty} \frac{n^\alpha}{(1+p)^n}=0
		$$
	\end{enumerate}
	
\end{proposition}
\begin{proof}
	（1）$\forall \varepsilon>0~$ ,考虑 $~\left|\dfrac{1}{(n)^{\alpha}}\right|<\varepsilon$
	取 $N=\left[\left(\dfrac{1}{\varepsilon}\right)^{1 / \alpha}\right]+1 $,$~ \forall n \rightarrow N$ 即可
	
	(2) $a>1 ~$ 时, 令$~y_n= \sqrt[n]{a}-1~,$那么 
	
	$$
	\begin{aligned}
		a=\left(1+y_n\right)^n&=1+n y_n+\cdots \\
		&>1+n y_n
	\end{aligned}
	$$
	
	
	也即
	
	$$
	y_n=\sqrt[n]{a}-1<\frac{a-1}{n}
	$$
	
	$\forall \varepsilon>0$ ，考虑 $$|\sqrt[n]{a}-1|<\frac{a-1}{n}<\varepsilon$$
	取 $~N=\left[\dfrac{a-1}{\varepsilon}\right]+10 ~$，则$~ \forall n >N$ 时
	
	$$
	|\sqrt[n]{a}-1|<\varepsilon
	$$
	
	$0<a<1$ 时 $\left(\dfrac{1}{a}\right)^{1 / n} \rightarrow 1$ 故 $(a)^{1 / n} \rightarrow 1$
	
	$a=1$ 时 ,显然
	
	（3）令$\sqrt[n]{n}-1=y_n$, 那么
	
	$$
	n=\left(1+y_n\right)^n=1+n y_n+\dfrac{n(n-1)}{2} y_n^2+\cdots 
	>\dfrac{n(n-1)}{2} y_n^2$$
	
	
	也即
	
	$$
	y_n<\sqrt{\dfrac{2}{n-1}}<\varepsilon
	$$
	
	
	取 $N=\left[\dfrac{2}{\varepsilon^2}+1\right]+10, ~~ \forall n>N$ 时
	
	$$
	|\sqrt[n]{n}-1|<\varepsilon
	$$
	
	
	对 $\lim\limits_{n \rightarrow \infty} \sqrt[n]{n^k}=1$ ，利用极限的乘法即可
	
	（4）设 $k \in \mathbb{Z}_{+}, ~ k > \alpha$ ，当 $n>k$ 时
	
	$$
	(1+p)^n>C_n^k p^k>\dfrac{n^k p^k}{k!} .
	$$
	
	
	从而 $\quad 0<\dfrac{n^\alpha}{(1+p)^n}<\dfrac{k!}{p^k} n^{\alpha-k}$
	且 $\alpha-k<0 ~$ 由（1）$~ n^{\alpha-k} \rightarrow 0$
	
	由夹逼性 $$\quad \lim _{n \rightarrow \infty} \dfrac{n^\alpha}{(1+p)^n}=0$$
	
\end{proof}
\begin{example}
	已知 $\lim\limits_{n \to \infty} a_n = a$，证明：
	$$\lim\limits_{n\to\infty}\dfrac{1}{2^n}\sum_{k=0}^{n}{n \choose k}a_k=a$$
	
\end{example}
\begin{proof}
	
	$$ \left| \dfrac{1}{2^n} \sum_{k=0}^{n} {n \choose k} a_k - a \right| = \left| \dfrac{1}{2^n} \sum_{k=0}^{n} {n \choose k} (a_k - a) \right| \leqslant \dfrac{1}{2^n} \sum_{k=0}^{n} {n \choose k} |a_k - a| $$
	$\forall \varepsilon > 0, \exists N_1 \text{ 使得当 } k > N_1 \text{ 时 }, |a_k - a| < \varepsilon$, 对所有的 $k$ 成立.
	
	\begin{align*}
		\frac{1}{2^n} \sum_{k=0}^{n} {n \choose k} |a_k - a| &= \frac{1}{2^n} \sum_{k=0}^{N_1} {n \choose k} |a_k - a| + \frac{1}{2^n} \sum_{k=N_1+1}^{n} {n \choose k} |a_k - a| \\
		&< \frac{1}{2^n} \sum_{k=0}^{N_1} {n \choose k} |a_k - a| + \varepsilon \\
		&< \frac{M (1 + n + \cdots + n^{N_1})}{2^n} + \varepsilon
	\end{align*}
	
	当 $n \to \infty$ 时, $\dfrac{1 + n + \cdots + n^{N_1}}{2^n} \to 0$，$M$ 是 $|a_k - a|$ 在 $N_1$ 内的上界.
	
	取 $N_2$,对所有 $n > N_2$,都有
	
	$$ \frac{M (1 + n + \cdots + n^{N_1})}{2^n} < \varepsilon $$
	
	取 $n > \max\{N_1, N_2\}$,就有
	
	$$ \left| \frac{1}{2^n} \sum_{k=0}^{n} {n \choose k} a_k - a \right| < 2\varepsilon $$
	
\end{proof}







\begin{proposition}(Cesaro平均定理)
	若$\lim\limits_{n\to\infty}a_n=a,$那么
	$$\lim\limits_{n\to\infty}\frac{a_1+\cdots+a_n}{n}=a$$
	\label{Cesaro}
\end{proposition}
\begin{proof}
	假设$a=0$,$\forall\varepsilon>0,\exists N_1\in \Z_{+},\forall n>N_1$,有$|a_n|<\varepsilon$
	$$\begin{aligned}
		\left|\frac{a_1+\cdots+a_n}{n}\right| & =\left|\frac{a_1+\cdots+a_{N_1}}{n}+\frac{a_{N_1+1}+\cdots+a_n}{n}\right| \\
		& \leqslant\left|\frac{a_1+\cdots+a_{N_1}}{n}\right|+\left|\frac{a_{N_1+1}+\cdots+a_n}{n}\right| \\
		& <\left|\frac{a_1+\cdots+a_{N_1}}{n}\right|+\varepsilon
	\end{aligned}$$
	对于固定的 $\varepsilon, N_1,\left|\dfrac{a_1+\cdots+a_{N_1}}{n}\right|$ 是固定的
	
	设 $c=a_1+\cdots+a_{N_1}$ ,想使 $\left|\dfrac{c}{n}\right|$ 与 $\varepsilon$ 有关,使 $\left|\dfrac{c}{n}\right|<\varepsilon$ ,取 $N_2=\left[\dfrac{c}{\varepsilon}\right]+100 n>N_2$ 时，$\left|\dfrac{c}{n}\right|<\varepsilon$
	
	$\forall n>\max \left[N_1, N_2\right] ~,~\left|\dfrac{a_1+\cdots+a_n}{n}\right|<\varepsilon+\varepsilon=2 \varepsilon$,即
	
	
	$$\lim _{n \rightarrow \infty} \frac{a_{1}+\cdots+a_n}{n}=a$$
	
\end{proof}
\begin{corollary}
	已知 $\lim\limits_{n \to \infty} a_n = a$，证明：
	
	$$ \lim_{n \to \infty} \sqrt[n]{a_1 \cdots a_n} = a $$
	\label{cesaro's co}
\end{corollary}
\begin{proof}
	设 $a_n > 0$ 且 $a > 0$（否则结论需要单独处理）. 由 $\lim\limits_{n \to \infty} a_n = a$ 和命题\ref{Cesaro}，有
	
	$$ \lim_{n \to \infty} \dfrac{\dfrac{1}{a_1} + \cdots + \dfrac{1}{a_n}}{n} = \dfrac{1}{a} $$
	
	因此
	
	$$ \lim_{n \to \infty} \dfrac{n}{\dfrac{1}{a_1} + \cdots + \dfrac{1}{a_n}} = a $$
	
	由基本不等式
	
	$$ \lim_{n \to \infty} \sqrt[n]{a_1 \cdots a_n} = a $$
	
	
\end{proof}
命题\ref{Cesaro}与命题\ref{cesaro's co}在下一节中将使用Stolz定理证明.

\begin{example}
	
	$a_i \geqslant 0$，证明：
	
	$$ \lim_{n \to \infty} (a_1^n + \cdots + a_p^n)^{1/n} = \max \{a_i\} $$
	
	
	
\end{example}
\begin{proof}
	由算术几何平均不等式,
	
	$$ \max \{a_i\} \leqslant (a_1^n + \cdots + a_p^n)^{1/n} \leqslant \max \{a_i\} \cdot p^{1/n} $$
	
	对于 $p \geqslant 1$,由夹逼性,得证.
\end{proof}
\begin{example}
	关于极限定义的一些反例:
	
	(1) $\forall \varepsilon > 0, \exists N$, 使得 $\forall n > N$，$x_n < \varepsilon$.
	
	取 $x_n = -n$，$N = 1$，$x_1, \ldots, x_n, \ldots < 0 < \varepsilon$.
	
	(2) $\forall \varepsilon > 0$，无穷多 $x_n$，$x_n < \varepsilon$.
	
	取 $x_n = (-1)^n$.
\end{example}

\begin{proposition}
	$\forall k \in \mathbb{Z}^+$，
	
	$$ \lim_{n \to \infty} x_n = a \Leftrightarrow \lim_{n \to \infty} x_{n+k} = a $$
\end{proposition}
\begin{proof}
	$\Rightarrow$: $\forall \varepsilon > 0, \exists N, \forall n > N$, 有 $|x_n - a| < \varepsilon$.
	
	对于 $n + k > N$，有 $|x_{n+k} - a| < \varepsilon$.
	
	因此，
	$$ \lim_{n \to \infty} x_{n+k} = a $$
	
	$\Leftarrow$: $\forall \varepsilon > 0, \exists N, \forall n + k > N$, 有 $|x_{n+k} - a| < \varepsilon$.
	
	对于 $n > N + k$，有 $|x_n - a| < \varepsilon$.
	
	因此，
	
	$$ \lim_{n \to \infty} x_n = a $$
\end{proof}
\begin{example}
	对于任意 $|q| > 0$，记明 $\left\{q^n\right\}$ 是无穷大量.
	
\end{example}
\begin{proof}
	$\forall G> 1$，考虑 $q^n > G$.
	
	取 $N = \left\lfloor \dfrac{\lg G}{\lg q} \right\rfloor + 1$，则对于 $\forall n > N$，成立
	$$ q^n > G$$
\end{proof}
\begin{example}
	$x_n \neq 0$，$\{x_n\}$ 是无穷大量当且仅当 $\{\dfrac{1}{x_n}\}$ 是无穷小量
\end{example}
\begin{proof}
	若 $\{x_n\}$ 是无穷大量，取 $G = \dfrac{1}{\varepsilon}$，
	
	$\forall G > 0, \exists N, \forall n > N$ 时，$|x_n| > G= \dfrac{1}{\varepsilon} \Leftrightarrow |x_n| > \dfrac{1}{\varepsilon} \Leftrightarrow \left|\dfrac{1}{x_n}\right| < \varepsilon$，即 $\left\{\dfrac{1}{x_n}\right\}$ 是无穷小量.
	
	若 $\left\{\dfrac{1}{x_n}\right\}$ 是无穷小量，取 $\varepsilon = \dfrac{1}{G}$，
	
	$\forall G > 0, \exists N, \forall n > N$ 时，$\left|\dfrac{1}{x_n}\right| < \varepsilon = \dfrac{1}{G} \Leftrightarrow |x_n| > G$，即 $\{x_n\}$ 是无穷大量.
\end{proof}


